% !TEX root = ../../../tree_notes.tex
% =====================================================================
%  11.3  The Binary (XOR) Trie                  sec:c08-binary
%  Exports: def:c08-bittrie  lem:c08-greedybit
%           fig:c08-bittriefig  lst:c08-maxxor
%           prob:c08-maxxor  lst:c08-maxxor-soln  ex:c08-xor-run   (practice: CF 1285D)
% =====================================================================
\section{The Binary (XOR) Trie}\label{sec:c08-binary}

\begin{previouslybox}
\begin{enumerate}[nosep, leftmargin=1.6em]
  \item In \secref{c08-triebasics} a node had up to $|\Sigma|$ children, one per letter. What are the
        ``letters'' of a non-negative integer written in binary? (the bits $0$ and $1$; the alphabet is
        $\{0,1\}$)
  \item XOR a bit with its \emph{opposite} and you get \emph{what}? ($1$ -- so to make a $1$ in the answer
        at a bit, the two operands must differ there)
  \item When comparing two binary numbers, which bit dominates the magnitude -- the high one or the low
        one? (the most significant set bit -- so greedy-from-the-top is safe)
  \item A max-heap or BST finds the extreme of a set \emph{by value}; a bitwise trie walks keys by their
        \emph{binary expansion}. What question about integers does that layout answer fastest? (an XOR
        extremum -- \chref{c05} gives order, not bit-wise complement)
\end{enumerate}
\end{previouslybox}

\begin{topicintro}
Change one thing and the whole machine reboots: instead of letters, read a key \emph{one bit at a time,
most significant first}. A \term{binary trie} (bitwise trie) is the \secref{c08-triebasics} structure over
$\Sigma=\{0,1\}$ -- every node has at most two children, and an integer $<2^{31}$ becomes a $31$-level
root-to-leaf path (\defref{c08-bittrie}). The point is not membership; it is \textbf{XOR optimisation}.
XOR sets bit $b$ to $1$ exactly when the two operands \emph{differ} at $b$, and higher bits outweigh all
lower ones combined -- so to build the \emph{largest} possible $x\oplus y$, stand at the top and, at each
level, greedily take the child holding the \emph{opposite} bit of $x$'s, whenever such a child exists
(\lemref{c08-greedybit}). That single greedy is why the greedy-from-MSB is provably optimal: you are
securing the highest-weight bit first and never trade it away.

The running example is $A=[3,10,5,25,2,8]$, the official LeetCode~421 input, whose answer is $28$. We
insert each value's bits, then query each value against the trie, preferring the complement bit; querying
$5$ steers toward $25$ and reads off $11100_2=28$ (\exref{c08-xor-run}). This is the same ``walk a path,
pick the child you need'' you already know from \secref{c08-triebasics}; the only new idea is the
\emph{opposite-bit} rule and \emph{why top-down} it is optimal. We also flag the array-index-vs-BST
contrast and hand you the divide-and-conquer-on-bits alternative (the CF~1285D practice).
\textbf{Prerequisites.} \secref{c08-triebasics} verbatim; bit operators $\oplus$, shift, mask; the
observation that a number's magnitude is set by its highest bit. No heap here (contrast \chref{c07}); the
binary trie is a different extreme-extraction -- extreme of \emph{XOR}, not of \emph{value}.

\textbf{After this section you can}
\begin{itemize}[nosep]
  \item insert integers as MSB-first bit paths and query max-XOR against a set in $\bigO{31}$; \hfill\lv{4}
  \item prove the opposite-bit greedy optimal from the dominance of the high bit; \hfill\lv{4}
  \item recognise a ``maximise XOR'' problem and reach for the bitwise trie over $\bigO{n^2}$ pairing;
        \hfill\lv{4}
\end{itemize}
\end{topicintro}

\begin{recognitionbox}[Pattern card -- maximise / minimise an XOR]
\reckey{the objective contains $x \oplus y$ and asks for its maximum (or a minimum under a constraint),
or asks to pair a query value against a set to make the XOR as large as possible -- the value $\le 10^9$
so its bit-length ($30$/$31$) is the working dimension.}
\recuse{
  \item ``maximum XOR of two numbers in an array'' -- insert all, query each (\probref{c08-maxxor});
  \item ``maximum XOR of $x$ with some array element'' (queries given) -- build once, query per $x$;
  \item ``maximize XOR subject to $a_j \le m$ / within a range / a past version'' -- the same walk, but
        the branches are filtered by counts or versions (\secref{c08-variants}, \probref{c08-persistxor}).
}
\recnotuse{
  \item the objective is \emph{sum} or \emph{product} or plain value extrema -- use a heap (\chref{c07})
        or a scan, not a bit path;
  \item you must find an XOR \emph{target} among pairs (``is there a pair with $a_i\oplus a_j = t$'') --
        a hash set of complements ($t\oplus a_i$) is $\bigO{n}$; a trie is overkill unless you need the
        \emph{max/min};
  \item ``minimise the \emph{maximum} $a_i \oplus X$'' over a single $X$ -- that is the D\&C-on-bits
        recursion (CF~1285D practice), a sibling of this greedy, not the pairwise query.
}
One-line test: \emph{is there an $\oplus$ whose value I am pushing to an extreme, and can I decide it
bit by bit from the top down? Then the binary trie greedy is the tool.}
\end{recognitionbox}

\subsection{Intuition}

Write $A=[3,10,5,25,2,8]$ in binary (using $5$ bits suffices since the largest, $25$, is $11001_2$):
\[
3{=}00011,\quad 10{=}01010,\quad 5{=}00101,\quad 25{=}11001,\quad 2{=}00010,\quad 8{=}01000.
\]
Now picture a trie whose edges are labelled $0$/$1$ and whose levels are these bit positions from
\textbf{left to right} (MSB first). Every number is a root-to-depth-$5$ path (\figref{c08-bittriefig}).
To maximise $5 \oplus y$, I want the answer's bits to be $1$ from the top: at bit~$4$ of $5$ (which is
$0$) I \emph{want} a $y$ whose bit~$4$ is $1$ -- and $25$ is the only one. Once I step onto the
$1$-branch, all of $y$'s lower bits are locked to $25$'s, but that is exactly what I want: the top bit
alone contributes $16$, more than every lower bit combined, so I never regret spending the choice to grab
it. Following complements down, I land on $25$ and read $5\oplus 25 = 11100_2 = 28$ -- the maximum over
all pairs.

That is the entire algorithm: the greedy is forced by bit \emph{weights}. A BST ordered by value
(\chref{c05}) cannot express this, because the partner maximising an XOR is neither the predecessor nor
the successor of $x$ -- $25$ is neither, yet it wins. The trie orders keys \emph{by their binary spelling},
and that is the right axis for $\oplus$.

\subsection{Definitions and Notation}

\begin{definition}{Binary (bitwise) trie}{c08-bittrie}
A \term{binary trie} over non-negative integers stores each value as a fixed-depth path of its bits,
\textbf{most significant first}: from the root, at level $b$ (bit $b$ of the value) take the $0$-child or
the $1$-child according to that bit. With values $<2^{B}$ the tree has depth $B$ (we use $B=31$ for
$<2^{31}$) and every node has at most two children. No terminal flag is needed for pure membership -- a
value ends at depth $B$ -- but a per-node \emph{count} $\mathit{cnt}$ (numbers passing through) is added
the moment we must support deletion or ``is this branch non-empty'' tests (\probref{c08-summit}).
\end{definition}

\noindent\textbf{In plain words.} Bits replace letters, MSB replaces ``first character,'' and depth is the
bit-width rather than key length. Everything else -- walk down, one node per bit -- is
\secref{c08-triebasics} unchanged.

\begin{lemma}{Opposite-bit greedy from the top is optimal}{c08-greedybit}
To maximise $x \oplus y$ over a stored set, process bits from MSB to LSB; at bit $b$, if a number in the
current candidate set has bit $b$ equal to $\lnot x_b$ (the complement), go to that child and set answer
bit $b$ to $1$; otherwise you must go to the same-bit child (answer bit $b$ is $0$). The resulting value
is the maximum achievable $x\oplus y$ over the set.
\end{lemma}
\begin{proofbox}
Induct on bit position from the top. Bit $b$ has weight $2^b$, and $\sum_{k<b} 2^k = 2^{b}-1 < 2^b$, so
\emph{any} number with answer-bit $b$ set exceeds \emph{every} number with it unset, regardless of lower
bits. Hence a maximum must set the highest possible bit first: at bit $b$ it is optimal to secure a $1$
whenever some candidate allows it (the complement branch is non-empty), because no later choice can
compensate for leaving it $0$. Among candidates sharing the chosen top bits, the same argument recurses
to lower bits. The trie's child at bit $b$ under the current node is exactly the set of candidates
matching the prefix already fixed, so ``go complement if it exists'' realises the induction step; the path
followed ends at a real stored value (every trie branch is a full-length value), so the greedy's partner
exists. \end{proofbox}

\begin{insightbox}
The greedy works \emph{because} XOR bit weights are powers of two and thus ``higher bit dominates'' --
a property an arbitrary objective lacks. This is the same reason radix sort reads from the top and a
binary search picks by the leading bit. When the constraint is richer than ``pick any stored $y$'' -- e.g.\
$y \le m$ or ``$y$ from the first $k$ insertions'' -- the pure greedy must be \emph{told} which branches
are still legal: a per-node count does it (\probref{c08-summit}), and a version-rooted count does the
range/version case (\secref{c08-variants}). The greedy skeleton never changes; only the ``is this child
available?'' predicate grows.
\end{insightbox}

\subsection{The Algorithm: Insert Bits, Query Complements}

We fix bit-width $B$ (here $31$, indices $30$ down to $0$). \code{insert(x)} walks $x$'s bits, creating
missing children, and bumps $\mathit{cnt}$ on each node (cheap insurance for deletion / non-empty tests).
\code{queryMaxXor(x)} walks again, at each bit choosing the complement child when present and setting that
answer bit, else the same child.

\begin{lstlisting}[caption={Binary trie for max XOR: \code{insert} sets bits MSB-first, \code{queryMaxXor} prefers the complement}, label={lst:c08-maxxor}]
struct Node { Node* ch[2]{}; int cnt = 0; };        // ch[0], ch[1]; cnt = how many values pass through
static const int MAXB = 30;                          // values < 2^31 -> bits 30..0
void insert(Node* root, int x) {
    Node* u = root;
    for (int b = MAXB; b >= 0; --b) {
        int bit = (x >> b) & 1;                      // this level's symbol is the bit
        if (!u->ch[bit]) u->ch[bit] = new Node();
        u = u->ch[bit];
        u->cnt++;                                     // maintain non-empty info for deletions
    }
}
int queryMaxXor(Node* root, int x) {                 // max x XOR y over y currently in the trie
    Node* u = root; int best = 0;
    for (int b = MAXB; b >= 0 && u; --b) {
        int bit = (x >> b) & 1, want = bit ^ 1;       // prefer the OPPOSITE bit
        if (u->ch[want] && u->ch[want]->cnt > 0) { best |= (1 << b); u = u->ch[want]; }
        else u = u->ch[bit];                           // else forced onto the same-bit child
    }
    return best;                                       // best is the max x XOR value, not the y
}
\end{lstlisting}

\begin{complexitybox}
Time: $\bigO{B}$ per insert and per query with $B=31$, i.e.\ constant $\bigO{31}$ -- exactly
\thmref{c08-searchcost} with $L=B$. Finding the max pair is $\bigO{nB}$: insert all $n$ (or insert
progressively), then query each. Memory: at most $n\cdot B$ nodes worst case, but shared prefixes make it
far smaller -- and a node is only $2$ pointers, so the memory pressure that plagued \code{child[26]}
(\defref{c08-nodelayout}) mostly vanishes at $\Sigma=\{0,1\}$. Verdict: the binary trie turns a
$\bigO{n^2}$ pairwise max into $\bigO{nB}$.
\end{complexitybox}

\subsection{Worked Example}

\begin{example}{Max-XOR query for $5$ against $A=\{3,10,5,25,2,8\}$ (using bits $4\ldots 0$)}{c08-xor-run}
Represent $A$ in $5$ bits (as above). Query $x=5=00101$:
\begin{itemize}[nosep]
  \item bit $4$ of $x$ is $0$; want $1$. Present (only $25$ has $1$). Go $1$; $\mathit{ans}$ bit $4=1$.
  \item bit $3$ of $x$ is $0$; want $1$. On the $1$-branch (holding $25{=}1\,1001$) next bit is $1$. Go $1$;
        bit $3=1$.
  \item bit $2$ of $x$ is $1$; want $0$. $25$'s bit $2$ is $0$. Go $0$; bit $2=1$.
  \item bit $1$ of $x$ is $0$; want $1$. $25$'s bit $1$ is $0$; no complement -> forced to $0$; bit $1=0$.
  \item bit $0$ of $x$ is $1$; want $0$. $25$'s bit $0$ is $1$; no complement -> forced to $1$; bit $0=0$.
\end{itemize}
$\mathit{ans}=11100_2=28$, partner $y=25$. Since $5\oplus 25=28$ and enumerating all $\binom{6}{2}$ pairs
gives $28$ as the maximum, the greedy is confirmed on the official answer.
\end{example}
\begin{tracebox}[Greedy complement choices for each query of $A$; running max of $x\oplus y$]
insert all of A, then query each x, taking the complement child whenever its cnt > 0.\\
$x{=}3\to$ partner $25$, $3\oplus 25=26$.\quad $x{=}10\to 25$, $10\oplus 25=19$.\quad $x{=}5\to 25$, $5\oplus 25=28$.\\
$x{=}25\to 5$, $25\oplus 5=28$.\quad $x{=}2\to 25$, $2\oplus 25=27$.\quad $x{=}8\to 25$, $8\oplus 25=17$.\\
Running max over queries $=\max(26,19,28,28,27,17)=28$. (The pair $\{5,25\}$ is hit twice -- querying $5$
finds $25$ and querying $25$ finds $5$ -- so the global maximum $28$ is discovered.)
\end{tracebox}

\begin{figure}[H]
\centering
\begin{tikzpicture}[diagram, scale=0.82, transform shape,
  level distance=8mm,
  level 1/.style={sibling distance=30mm}, level 2/.style={sibling distance=15mm},
  level 3/.style={sibling distance=10mm}, level 4/.style={sibling distance=8mm},
  level 5/.style={sibling distance=6mm},
  bit/.style={circle, draw=darkgray, fill=white, minimum size=4.5mm, font=\scriptsize},
  tend/.style={circle, draw=successgreen, fill=successgreen!20, minimum size=4.5mm, font=\scriptsize}]
  \node[bit] (rt) {$\cdot$}
    child { node[bit] (z0) {0}
      child { node[bit] (z1) {0}
        child { node[bit] (z2a) {0}
          child { node[bit] (z3a) {1} child { node[tend] (e3) {1} } } }
        child { node[bit] (z2b) {1}
          child { node[bit] (z3b) {0} child { node[tend] (e5) {1} } } } } }
    child { node[bit] (o0) {1}
      child { node[bit] (o1) {1}
        child { node[bit] (o2) {0}
          child { node[bit] (o3) {0}
            child { node[tend] (e25) {1} } } } } };
  \draw[tedgehl] (rt) -- (o0) -- (o1) -- (o2) -- (o3) -- (e25);
  \node[annot, align=left, font=\scriptsize] at (-4.6,-1.4)
    {left spine $00011=3$,\\ $00101=5$;\\ right spine $11001=25$};
  \node[annot, align=left, font=\scriptsize, text=warningred] at (4.9,-2.6)
    {\textbf{greedy for }$x{=}5$:\\ at every level jump\\ to the complement\\ child -- lands on $25$\\
    ($5\oplus 25=28$)};
\end{tikzpicture}
\caption{The binary trie of $\{3,5,25\}\subset A$ (three of the six values, $5$ bits, MSB at the top).
Each node prints the bit taken at that level; the highlighted red root-to-leaf walk is the opposite-bit
greedy for querying $x=5$, ending at $25$ and yielding $28$. The full trie of $A$ adds the branches for
$10{=}01010$, $2{=}00010$, $8{=}01000$; the production code uses $31$ levels (\coderef{c08-maxxor}).}\label{fig:c08-bittriefig}
\end{figure}

\subsection{Variations and Twists}

\begin{itemize}[nosep]
  \item \textbf{Pairwise vs query form}: max \emph{pair} in one array (LC~421) -- insert then query each.
        Max $x\oplus a_i$ for \emph{given} $x$'s -- build once, query each $x$. Same code, loop inverted. \hfill\lv{4}
  \item \textbf{Minimum XOR pair}: order the values, then the min XOR pair is adjacent after sorting
        (a totally different, cheaper route); the trie is for \emph{max}. \hfill\lv{4}
  \item \textbf{Constrained max}: ``$y$ within $[l,r]$ of a value or an index, or a past version'' needs
        the branch-availability predicate upgraded from ``child exists'' to ``child has count/roots in
        range'' -- precisely \secref{c08-variants} and \probref{c08-persistxor}. \hfill\lv{5}
  \item \textbf{D\&C on bits (no trie)}: for ``minimise the max $a_i\oplus X$'' (CF~1285D) you can split
        the array by each bit and recurse, saving the explicit structure -- a sibling technique. \hfill\lv{4}
\end{itemize}

\subsection{Common Pitfalls}

\begin{warningbox}[Trap 1 -- reading bits LSB-first]
The greedy \emph{requires} MSB-first (\lemref{c08-greedybit}): the top bit dominates. If you insert or
query from bit $0$ upward, you commit to low bits before high ones and the ``prefer complement'' rule
 maximises the wrong thing -- \textbf{WA}. Loop `for (b = MAXB; b >= 0; --b)` everywhere.
\end{warningbox}
\begin{warningbox}[Trap 2 -- wrong bit-width]
$B$ must cover the largest possible value: $a_i < 2^{31}$ needs bits $30..0$ ($B=31$). Using $30$ bits
(max bit index $29$) silently drops the top bit and misses pairs differing only there. When a constraint
says $\le 10^9 < 2^{30}$, index $29..0$ ($B=30$) suffices -- match $B$ to the bound, not habit.
\end{warningbox}
\begin{warningbox}[Trap 3 -- returning the partner instead of the XOR]
\code{queryMaxXor} returns $\mathit{best}$, the \emph{value of $x\oplus y$}, by setting answer bits. If you
return the traversed leaf's stored number instead you have $y$, not the objective. Keep a running
\code{best |= (1<<b)} as in \coderef{c08-maxxor}.
\end{warningbox}

\subsection{Problem Lab}

The section's anchor is LeetCode~421, \probref{c08-maxxor}: given an array, return the maximum value of
$a_i \oplus a_j$ over all pairs. It is the cleanest ``insert then query complements'' use of the bitwise
trie, and the template every constrained-XOR problem (LC~1707, the Summit CF~706D) extends.

% ---------------------------------------------------------------------------
\solvedproblem{c08-maxxor}{Maximum XOR of Two Numbers in an Array}{LeetCode 421}
  {https://leetcode.com/problems/maximum-xor-of-two-numbers-in-array/}{LC Medium · L4}
  {binary trie MSB-first, opposite-bit greedy, max pair, bit-width 31}

\begin{problemstatement}[Maximum XOR of Two Numbers in an Array]
Given an integer array \texttt{nums}, return the \textbf{maximum result} of \texttt{nums[i] XOR nums[j]}
over all pairs $0 \le i, j < n$. Constraints (official): $1 \le n \le 2\cdot 10^{5}$ and
$0 \le \texttt{nums[i]} \le 2^{31} - 1$.
\end{problemstatement}

\begin{attemptgate}[Attempt it yourself first -- \lv{4}, give yourself 20 minutes]
$\bigO{n^2}$ pairs is too slow for $n=2\cdot 10^5$. Reframe: for each value $x$, find the stored $y$ that
maximises $x\oplus y$ -- that is the trie greedy. Insert everything first, then query each $x$.
\peek{Reading bits top-down, which child maximises the answer's current bit?}
\nudge{The child whose bit is the \emph{opposite} of $x$'s current bit, if it exists -- it sets that answer
bit to $1$. Walk MSB to LSB, greedily preferring the complement branch, accumulating the answer.}
\fullhint{Build a $31$-level binary trie of all values (\coderef{c08-maxxor}). For each $x$, descend:
at bit $b$, if $x$'s complement child exists, go there and set answer bit $b$; else go to the same-bit
child. Take the max over all queries. $O(nB)$ total.}
\end{attemptgate}

\solstep{Understanding the Problem}
We need the largest XOR among all pairs. Brute force is $\binom{n}{2}$ pairs -- $O(n^2)$, too slow at
$n=2\cdot 10^5$. The objective is an \emph{XOR extremum}, which is exactly what a bitwise trie answers:
fix a value $x$ and grab the $y$ that differs from it in the highest bits. Doing that for every $x$ and
taking the best discovers the global max pair (the true pair contributes when you query either of its
members). Bit-width must reach $2^{31}$, so we use levels $30..0$.

\solstep{Key Observations}
\begin{observation}
Maximising $x\oplus y$ is a per-bit greedy from the MSB down, valid because $2^b > \sum_{k<b}2^k$
(\lemref{c08-greedybit}) -- secure the high bit first, never trade it.
\end{observation}
\begin{observation}
A binary trie answers ``is there a stored value with this bit-prefix'' in one $0/1$ pointer hop, so each
query is $\bigO{31}$ (\defref{c08-bittrie}).
\end{observation}
\begin{observation}
The maximum \emph{pair} is found by querying each element against the set of all elements: the winning
pair $(p,q)$ is discovered when we query $p$ (finding $q$) and again when we query $q$.
\end{observation}

\solstep{Think It Through With Me}
\begin{thinkbox}
\deadend{Enumerate all pairs.}
{$O(n^2)=4\cdot 10^{10}$ at $n=2\cdot 10^5$ -- \re. Ignores the bit structure.}
\deadend{Sort and check adjacent pairs.}
{Adjacency after sorting helps \emph{minimum} XOR, not maximum; the max-XOR partner of $x$ is neither
neighbour in value order (\figref{c08-bittriefig}: $25$ beats $5$ though neither is $5$'s successor).}
\aha{Insert every value's bits into a $31$-level trie, then for each $x$ walk it preferring the complement
bit at each level; accumulate and maximise. That is $\bigO{nB}$.}
\end{thinkbox}

\solstep{Solution}
Build a binary trie and insert all $n$ values (each an $\bigO{B}$ walk). For every value $x$, run
\code{queryMaxXor(x)} -- descend from bit $30$ to $0$, at each level taking the complement child when it
exists (setting that answer bit) else the same-bit child. Track the maximum returned value over all $x$
and output it. Because the true max pair is hit when querying either of its elements, taking the max over
all queries is the global answer. \coderef{c08-maxxor-soln} packages the two loops.

\solstep{Proof of Correctness}
\begin{proofbox}[The query returns the true max over pairs]
By \lemref{c08-greedybit}, for a fixed $x$ the opposite-bit descent returns $\max_{y\in A} x\oplus y$.
Let $(p,q)$ be a pair attaining the global maximum $M=\max_{i,j} a_i\oplus a_j$. When $x=p$, the query
returns $\max_y p\oplus y \ge p\oplus q = M$; and it cannot exceed $M$ since $M$ is the largest pair XOR
by definition, so it equals $M$. Thus the maximum over queries is $M$. Conversely every query value is
$p\oplus y$ for real stored $y$, hence $\le M$, so the returned maximum is exactly $M$. \end{proofbox}

\solstep{Implementation}
\begin{lstlisting}[caption={LeetCode 421 -- insert all values, then max-query each against the binary trie}, label={lst:c08-maxxor-soln}]
class Solution {
    struct Node { Node* ch[2]{}; };
    Node* rt = new Node();
    void insert(int x){ Node* u=rt; for(int b=30;b>=0;--b){int k=(x>>b)&1; if(!u->ch[k])u->ch[k]=new Node(); u=u->ch[k];} }
    int maxXor(int x){ Node* u=rt; int ans=0;
        for(int b=30;b>=0 && u;--b){ int k=(x>>b)&1, want=k^1;
            if(u->ch[want]){ ans|=(1<<b); u=u->ch[want]; } else u=u->ch[k]; }
        return ans; }
public:
    int findMaximumXOR(vector<int>& nums) {
        for (int x : nums) insert(x);          // stage 1: build trie of all values
        int best = 0;
        for (int x : nums) best = max(best, maxXor(x));   // stage 2: query complements
        return best;
    }
};
\end{lstlisting}

\solstep{Trace}
\begin{tracebox}[Official sample: nums = (3,10,5,25,2,8) -> 28]
After inserting all six (bits 30..0, only low bits non-zero):\\
query 3 : complement descent lands on 25 -> $3\oplus 25 = 26$.\\
query 10: descends to 25 -> $10\oplus 25 = 19$.\\
query 5 : descends to 25 -> $5\oplus 25 = 28$.\\
query 25: descends to 5 -> $25\oplus 5 = 28$ (the max, not $25\oplus 2 = 27$).\\
query 2 : descends to 25 -> $27$.\quad query 8: -> $17$.\\
Running max over queries $=\max(26,19,28,28,27,17)=28$. \checkmark\ matches official output.
\end{tracebox}

\solstep{Complexity}
\begin{complexitybox}
Time: $O(nB)$ with $B=31$: $n$ inserts and $n$ queries, each a $\bigO{31}$ walk -- about
$6\cdot 10^{6}$ operations at $n=2\cdot 10^5$, comfortably fast (vs $O(n^2)$ brute -- \re).
Memory: $O(nB)$ nodes worst case, typically far fewer from shared prefixes; each node is $2$ pointers, so
the $\Sigma=\{0,1\}$ layout is light (\defref{c08-nodelayout}). Verdict: \ac.
\end{complexitybox}

\solstep{Pitfalls and Variants}
\begin{warningbox}
Live slips: (i) LSB-first loop (Trap 1) -- \textbf{WA}; (ii) bit-width $29..0$ for a $2^{31}$ bound
(Trap 2) -- misses the top bit; (iii) returning the leaf value instead of the accumulated XOR (Trap 3).
\emph{Variant:} a slicker $O(nB)$ alternative (LC~421 editorial ``Solution 2'') greedily builds a target
prefix using a hash set of masked prefixes and no explicit trie -- worth seeing, but the trie is the idea
that generalises to \secref{c08-variants} and the Summit.
\end{warningbox}

\begin{insightbox}[Key takeaways]
\begin{enumerate}[nosep, leftmargin=1.6em]
  \item A bit is a letter from $\Sigma=\{0,1\}$; the string trie of \secref{c08-triebasics} becomes a
        \emph{binary trie} the instant you read keys MSB-first (\defref{c08-bittrie}).
  \item XOR is maximised by the \emph{opposite-bit} greedy, and it is optimal because a high bit outweighs
        all low bits together (\lemref{c08-greedybit}).
  \item Max pair $=$ build trie of all, query complements for each -- $\bigO{nB}$ beats $\bigO{n^2}$
        (\probref{c08-maxxor}); the same walk with upgraded ``is child available?'' predicates handles
        constraints (\probref{c08-persistxor}, \probref{c08-summit}).
\end{enumerate}
\end{insightbox}

\begin{hardeningbox}
\begin{itemize}[nosep, leftmargin=1.4em]
  \item \textbf{Cold-write \coderef{c08-maxxor}} and reproduce \exref{c08-xor-run}: confirm querying $5$
        returns $28$ and the running max over $A$ is $28$.
  \item \textbf{Add deletion}: use the $\mathit{cnt}$ field to erase a value (decrement along its path) and
        make \code{queryMaxXor} skip children with $\mathit{cnt}=0$ -- this is the Summit's move
        (\probref{c08-summit}).
  \item \textbf{Read ahead}: to query ``max XOR with a $y\le m$ or from a version'' you need per-branch
        counts that survive into the past -- \secref{c08-variants} builds exactly that.
\end{itemize}
\end{hardeningbox}

\subsection{Practice Problems}

\begin{stuckbox}
\begin{itemize}[nosep, leftmargin=1.4em]
  \item Are you inserting \emph{and} querying from the \emph{same} top bit downward?
  \item Is your bit-width $B$ matched to the value bound, not hardcoded by habit?
  \item In the query, are you accumulating the \emph{XOR value} (answer bits), not the partner?
\end{itemize}
\end{stuckbox}

\begin{practiceset}
  \practice{Dr. Evil Underscores}{Codeforces 1285D}
    {https://codeforces.com/problemset/problem/1285/D}{CF 1600 · L4}
    {Choose a non-negative integer $X$ to \emph{minimise} $\max_i (a_i \oplus X)$, and output that minimum.
    Unlike \probref{c08-maxxor} this is not a pairwise query, but it is the same top-down bit view: a
    divide-and-conquer on bit position $b$ -- split the array into the group with bit $b$ zero and the
    group with bit $b$ one. If only one group is non-empty, $X$'s bit $b$ costs nothing and you recurse on
    $b-1$; if \emph{both} are non-empty, whatever $X$ does, one group gets a leading $1$ at $b$, so add
    $2^b$ plus the \emph{better} (min) of recursing on the two groups at $b-1$. The recursion is a trie
    built on the fly, so writing it \emph{as} a binary-trie traversal earns the same picture -- and answers
    the ``is a branch available?'' question you will need for \probref{c08-summit}.}
\end{practiceset}

\begin{recallbox}[10-second refresher -- 11.3 The Binary (XOR) Trie]
Read integers as MSB-first bit strings: a node has $\le 2$ children (\defref{c08-bittrie}). Maximise
$x\oplus y$ by walking top-down, \emph{preferring the complement bit} whenever that child exists
(\lemref{c08-greedybit}); it is optimal because a high bit outweighs all lower bits. Max pair $=$ insert
all, query each (\coderef{c08-maxxor}, LeetCode~421 \probref{c08-maxxor}, $O(nB)$ with $B=31$). A per-node
$\mathit{cnt}$ upgrades the walk to deletion and constrained queries (\probref{c08-summit},
\probref{c08-persistxor}). Mind the traps: MSB-first, correct bit-width, return the XOR not the partner.
\selftest{
  \item Why must bits be processed most-significant-first for the greedy?
  \item State \lemref{c08-greedybit} and the inequality that makes it correct.
  \item How does finding the max \emph{pair} reduce to $n$ max-XOR \emph{queries}?
  \item What single field turns the plain binary trie into one that supports deletion?
  \item Why is a BST (\chref{c05}) the wrong structure for maximum XOR?
}
\end{recallbox}
